\documentclass[a4paper,twoside]{article}
\usepackage{blindtext}  
\usepackage{geometry}

% Chinese support
\usepackage[UTF8, scheme = plain]{ctex}

% Page margin layout
\geometry{left=2.3cm,right=2cm,top=2.5cm,bottom=2.0cm}


\usepackage{listings}
\usepackage{xcolor}
\usepackage{geometry}
\usepackage{amsmath}
\usepackage{float}
\usepackage{hyperref}

\usepackage{graphics}
\usepackage{graphicx}
\usepackage{subfigure}
\usepackage{epsfig}
\usepackage{float}

\usepackage{algorithm}
\usepackage[noend]{algpseudocode}

\usepackage{booktabs}
\usepackage{threeparttable}
\usepackage{longtable}
\usepackage{listings}
\usepackage{tikz}

% cite package, to clean up citations in the main text. Do not remove.
\usepackage{cite}

\usepackage{color,xcolor}

%% The amssymb package provides various useful mathematical symbols
\usepackage{amssymb}
%% The amsthm package provides extended theorem environments
\usepackage{amsthm}
\usepackage{amsfonts}
\usepackage{enumerate}
\usepackage{enumitem}
\usepackage{listings}

\usepackage{indentfirst}
\setlength{\parindent}{2em} % Make two letter space in the first paragraph
\usepackage{setspace}
\linespread{1.5} % Line spacing setting
\usepackage{siunitx}
\setlength{\parskip}{0.5em} % Paragraph spacing setting

% \usepackage[contents =22920202204622, scale = 10, color = black, angle = 50, opacity = .10]{background}

\renewcommand{\figurename}{图}
\renewcommand{\lstlistingname}{代码} 
\renewcommand{\tablename}{表格}
\renewcommand{\contentsname}{目录}
\floatname{algorithm}{算法}

\graphicspath{ {images/} }

%%%%%%%%%%%%%
\newcommand{\StudentNumber}{22920202204622}  % Fill your student number here
\newcommand{\StudentName}{熊恪峥}  % Replace your name here
\newcommand{\PaperTitle}{作业（五）}  % Change your paper title here
\newcommand{\PaperType}{算法设计与分析} % Replace the type of your report here
\newcommand{\Date}{2022年3月23日}
\newcommand{\College}{信息学院}
\newcommand{\CourseName}{算法设计与分析}
%%%%%%%%%%%%%

%% Page header and footer setting
\usepackage{fancyhdr}
\usepackage{lastpage}
\pagestyle{fancy}
\fancyhf{}
% This requires the document to be twoside
\fancyhead[LO]{\texttt{\StudentName }}
\fancyhead[LE]{\texttt{\StudentNumber}}
\fancyhead[C]{\texttt{\PaperTitle }}
\fancyhead[R]{\texttt{第{\thepage}页，共\pageref*{LastPage}页}}


\title{\PaperTitle}
\author{\StudentName}
\date{\Date}

\lstset{
	basicstyle          =   \sffamily,          % 基本代码风格
	keywordstyle        =   \bfseries,          % 关键字风格
	commentstyle        =   \rmfamily\itshape,  % 注释的风格，斜体
	stringstyle         =   \ttfamily,  % 字符串风格
	flexiblecolumns,                % 别问为什么，加上这个
	numbers             =   left,   % 行号的位置在左边
	showspaces          =   false,  % 是否显示空格，显示了有点乱，所以不现实了
	numberstyle         =   \zihao{-5}\ttfamily,    % 行号的样式，小五号，tt等宽字体
	showstringspaces    =   false,
	captionpos          =   t,      % 这段代码的名字所呈现的位置，t指的是top上面
	frame               =   lrtb,   % 显示边框
}

\lstdefinestyle{PythonStyle}{
	language        =   Python, % 语言选Python
	basicstyle      =   \zihao{-5}\ttfamily,
	numberstyle     =   \zihao{-5}\ttfamily,
	keywordstyle    =   \color{blue},
	keywordstyle    =   [2] \color{teal},
	stringstyle     =   \color{magenta},
	commentstyle    =   \color{red}\ttfamily,
	breaklines      =   true,   % 自动换行，建议不要写太长的行
	columns         =   fixed,  % 如果不加这一句，字间距就不固定，很丑，必须加
	basewidth       =   0.5em,
}

\algnewcommand\algorithmicinput{\textbf{Input:}}
\algnewcommand\algorithmicoutput{\textbf{Output:}}
\algnewcommand\Input{\item[\algorithmicinput]}%
\algnewcommand\Output{\item[\algorithmicoutput]}%

\begin{document}
	
%%%%%%%%%%%%%%%%%%%%%%%%%%%%%%%%%%%%%%%%%%%%
\makeatletter % change default title style
\renewcommand*\maketitle{%
	\begin{center} 
		\bfseries  % title 
		{\LARGE \@title \par}  % LARGE typesetting
		\vskip 1em  %  margin 1em
		{\global\let\author\@empty}  % no author information
		{\global\let\date\@empty}  % no date
		\thispagestyle{empty}   %  empty page style
	\end{center}%
	\setcounter{footnote}{0}%
}
\makeatother
%%%%%%%%%%%%%%%%%%%%%%%%%%%%%%%%%%%%%%%%%%%%
	
	
\thispagestyle{empty}

\vspace*{1cm}

\begin{figure}[h]
	\centering
	\includegraphics[width=4.0cm]{logo.png}
\end{figure}

\vspace*{1cm}

\begin{center}
	\Huge{\textbf{\PaperType}}
	
	\Large{\PaperTitle}
\end{center}

\vspace*{1cm}

\begin{table}[h]
	\centering	
	\begin{Large}
		\renewcommand{\arraystretch}{1.5}
		\begin{tabular}{p{3cm} p{5cm}<{\centering}}
			姓\qquad 名 & \StudentName  \\
			\hline
			学\qquad号 & \StudentNumber \\
			\hline
			日\qquad期 & \Date  \\
			\hline
			学\qquad院 & \College  \\
			\hline
			课程名称 & \CourseName  \\
			\hline
		\end{tabular}
	\end{Large}
\end{table}

\newpage

\title{
	\Large{\textcolor{black}{\PaperTitle}}
}
	
	
\maketitle
	
\tableofcontents
 
\newpage
\setcounter{page}{1}

\begin{spacing}{1.2}

\section{题6.4}

对于维度序列为$\{5,10,3,12,5,50,6\}$的矩阵链，可以按递推方程\eqref{eqn:mat}填表~\ref{tbl:matm}和表~\ref{tbl:mats}

\begin{equation}
\label{eqn:mat}
\begin{aligned}
m_{i,j}=
\begin{cases}
	0 \ \ &i=j \\
	\mathop{\min}_{i\le k<j}\{m_{i,k}+m_{k,j}+p_{i-1}p_jp_k\} \ \ &i\ne j
\end{cases}
\end{aligned}
\end{equation}

\begin{table}[htbp]
	\label{tbl:matm}
	\centering
	\caption{$m_{i,j}$内容}
	\begin{tabular}{|c|c|c|c|c|c|}
		\hline
		0&150&330&405&1655&2010\\ \hline
		inf&0&360&330&2430&1950\\ \hline
		inf&inf&0&180&930&1770\\ \hline
		inf&inf&inf&0&3000&1860\\ \hline
		inf&inf&inf&inf&0&1500\\ \hline
		inf&inf&inf&inf&inf&0\\ \hline
	\end{tabular}
\end{table}

\begin{table}[htbp]
	\label{tbl:mats}
	\centering
	\caption{$s_{i,j}$内容}
	\begin{tabular}{|c|c|c|c|c|c|}
		\hline
		N&1&2&2&4&2\\\hline
		N&N&2&2&2&2\\\hline
		N&N&N&3&4&4\\\hline
		N&N&N&N&4&4\\\hline
		N&N&N&N&N&5\\\hline
		N&N&N&N&N&N\\\hline
	\end{tabular}
\end{table}

并且根据$s_{i,j}$中的内容可以得到括号加法为

$$
( ( A_{1} A_{2} ) ( ( A_{3} A_{4} ) ( A_{5} A_{6} ) ) ) 
$$

代码实现见\nameref{sec:app_impl}中的代码\ref{code:implmat}

\section{题6.6}

\begin{equation}
	\begin{aligned}
		P(n)=
		\begin{cases}
		\Theta(1) \ \ \ &n=1\\
		\mathop{\sum}_{k=1}^{n-1}P(k)P(n-k) \ \ \ &n\ge 2
		\end{cases}
	\end{aligned}
\end{equation}

要证$P(n)=\Omega (2^n)$只需证$P(n)\ge c\cdot 2^n$

\textbf{证：} $P(n)\ge c\cdot 2^n$

\textbf{当} $n=1$， $P(1)=1>c\cdot 2$，当$c\le 1$

\textbf{若} $\forall k<n, P(k) \ge c\cdot 2^k$成立，则

\begin{align*}
	P(n)&=\mathop{\sum}_{k=1}^{n-1}P(k)P(n-k)\\
	&\ge c\cdot 2^1 \times c\cdot 2^{n-1} + c\cdot 2^2 \times c\cdot 2^{n-2}+\cdots+  c\cdot 2^{n-1} \times  c\cdot2^1  \\
	&= c^2 \cdot (n-1) \cdot (2^n) \\
	&\ge c\cdot 2^n
\end{align*}

当

$$
c \ge \frac{1}{n-1}
$$

\section{题6.8}

由状态转移方程\eqref{eqn:assm}

\begin{equation}
\label{eqn:assm}
\begin{aligned}
f&=\min\{f_1[n]+x_1,f_2[n]+x_2\} \\
f_1[n]&=
\begin{cases}
	e_1+a_{1,1} \ \ \ &j=1 \\
	\min\{f_1[j-1],f_2[j-1]+t_{2,j-1}\}+a_{1,j} \ \ \ &j\ge2 
\end{cases} \\
f_2[n]&=
\begin{cases}
	e_2+a_{2,1} \ \ \ &j=1 \\
	\min\{f_2[j-1],f_1[j-1]+t_{1,j-1}\}+a_{2,j} \ \ \ &j\ge2 
\end{cases}
\end{aligned}
\end{equation}

可知$f_1[n]$的计算引用了$f_1[n-1]$和$f_2[n-1]$，$f_1[n-1]$的计算引用了$f_1[n-2]$和$f_2[n-2]$，\ldots
以此类推，设$f_1,f_2$中元素引用次数$r_1[n]=1, r_2[n]=1$，有

$$
r_1[n-1]=r_1[n]+r_2[n], r_2[n-1]=r_1[n]+r_2[n]
$$

易知

\begin{align*}
	r_1[n] &= O(2^n) \\
	r_2[n] &= O(2^n)
\end{align*}

则子问题都被多次计算，因此有重叠子问题的性质。

\section{题6.10}

根据状态转移方程可以实现算法~\ref{algo:lcsrecursive}

\begin{algorithm}
	\caption{LCS问题}
	\label{algo:lcsrecursive}
	\begin{algorithmic}[1]
		\Input{两个字符及其长度，以及备忘录的缓冲区}
		
		\Procedure{LCS}{$X$,$Y$,$m$,$n$,$memo$}     
		
		\If{$m==0 \mathop{or} n==0$}
		\State \Return $memo[i][j]=0$
		\EndIf
		
		\If{$memo[i][j]\ne -1$}
		\State \Return $memo[i][j]$
		\EndIf
		
		\If{$X[m]==Y[n]$}
		\State \Return $memo[i][j]=$ \Call{LCS}{$X$,$Y$,$m-1$,$n-1$,$memo$}$+1$
		\Else
		\State \Return  $memo[i][j]=\max \{$ \Call{LCS}{$X$,$Y$,$m-1$,$n$,$memo$}$,$\Call{LCS}{$X$,$Y$,$m$,$n-1$,$memo$}$\}$
		\EndIf
		
		\EndProcedure
	\end{algorithmic}
\end{algorithm}	

\section{题6.11}

$X_m$和$Y_n$是$A$开始的字符串，不妨设存在长度$\ge 2$的最长公共子序列，则

\begin{align*}
	X_m &= A\ldots s\ldots \\
	Y_n &= A\ldots s\ldots
\end{align*}

其中$s$可以是任意字符，表示起始的$A$之后任何一个重合的字符。

\textbf{若}存在一个最长公共子序列${lcs}_k$，${lcs}_k$不以$A$开头，则不妨设${lcs}_k$以$s$开头。

\textbf{则}可以构造${lcs}_k'$，令

$$
{lcs}_k'=A+{lcs}_k=As\ldots
$$

由于$X_m$和$Y_n$是$A$开始的字符串，按照定义显然${lcs}_k'$是$X_m$和$Y_n$的一个最长公共子序列。但

$$
\mathop{len}({lcs}_k')=\mathop{len}({lcs}_k)+1 > \mathop{len}({lcs}_k)
$$

与最长公共子序列的最优性矛盾，因此假设不成立，一切最长公共子序列都以$A$开头

\section{题6.16}
由\eqref{eqn:wbetter}计算$w$时每一次第二层循环需要1次或2次加法，因此相当于给时间复杂度贡献了常数。而

\begin{equation}
	\label{eqn:wbetter}
	\begin{aligned}
		w_{i,j}=
		\begin{cases}
			q_{i-1} \ \ \ \ &j=i-1 \\
			w_{i,j-1}+p_i+q_j &1\le i\le j\le n
		\end{cases}
	\end{aligned}
\end{equation}

使用\eqref{eqn:wworse}计算，则需要$j-i+1+j-i+2=2j-2i+3$次加法。

\begin{equation}
	w=\sum_{l=i}^jp_l+\sum_{l=i-1}^jq_l
\end{equation}

则相当于给时间复杂度贡献了一个乘积项$n$，时间复杂度会变成$\mathcal{O}(n^4)$

\section{题6.19}

首先注意到删除操作不是必要的。例如删除$b$使得$abcd$与$acd$相同，\emph{等效于}插入$b$使得$abcd$与$acd$相同。即在一个串中删除一个字符$c$等于在另一个串中对应位置插入一个$c$。则有如下三种不同的操作：

\begin{itemize}
	\item 在串$A$中插入
	\item 在串$B$中插入
	\item 替换一个字符。由于在$A$中和$B$中替换是等效的，不妨设在$A$中。
\end{itemize}

那么为了完成$A$中前$i$个字符和$B$中前$j$个字符的编辑，可以：

\begin{enumerate}
	\item 在完成了$A$中前$i$个字符和$B$中前$j-1$个字符的编辑之后，在$A$中插入$B[j]$
	\item 在完成了$A$中前$i-1$个字符和$B$中前$j$个字符的编辑之后，在$A$中插入$A[i]$
	\item 在完成了$A$中前$i-1$个字符和$B$中前$j-1$个字符的编辑之后，将$A[i]$和$B[j]$改写成一样的字符。
\end{enumerate}
特别地，对于情况3，当$A[i]=B[j]$时，可以不进行任何操作。对于空串$\varepsilon$，为了使它跟字符串$S$相同，可以插入$\mathop{len}(S)$次字符。

根据以上分析可以列出递推方程\eqref{eqn:dist}。

\begin{equation}
\label{eqn:dist}
\begin{aligned}
f_{i,j}=\begin{cases}
	i \ \ \ &j=0\\
	j \ \ \ &i=0\\
	\mathop{\min}\{f_{i,j-1},f_{i-1,j},f_{i-1,j-1}\}+1 \ \ \ &A[i-1]\ne B[j-1]\\
	\mathop{\min}\{f_{i,j-1}+1,f_{i-1,j}+1,f_{i-1,j-1}\} \ \ \ &A[i-1]=B[j-1]
\end{cases}
\end{aligned}
\end{equation}

据此实现算法\ref{algo:dist}。

\begin{algorithm}
	\caption{编辑距离问题问题}
	\label{algo:dist}
	\begin{algorithmic}[1]
		\Input{两个字符}
		
		\Procedure{distance}{$X$,$Y$,$m$,$n$,$memo$}     
		\For{$i=0\to |A|$}
		\State $f_{i,0}\gets i$
		\EndFor
		
		\For{$i=0\to |B|$}
		\State $f_{0,j}\gets j$
		\EndFor
		
		\For{$i=1\to |A|$}
		\For{$i=1\to |B|$}
		\If{$A[i-1]==B[i-1]$}
		\State $f_{i,j} \gets \mathop{\min}\{f_{i,j-1}+1,f_{i-1,j}+1,f_{i-1,j-1}\}$
		\Else
		\State $f_{i,j} \gets \mathop{\min}\{f_{i,j-1},f_{i-1,j},f_{i-1,j-1}\}+1$
		\EndIf
		\EndFor
		\EndFor
		
		\State \Return $f_{|A|,|B|}$
		\EndProcedure
	\end{algorithmic}
\end{algorithm}	

\newpage

\section{实验一}

要求给定两个程序，判断它们的相似性。显然，\emph{程序的相似性和代码字符串的相似性无关，而与实际执行逻辑的相似性有关}。那么最准确的方式是进行DFA(Data Flow Analysis)和CFA(Control Flow Analysis)，对于相似的程序它们应当能相当准确地反映出相似度。这正是现代IDE对重复代码给出修改建议的方式。但这种方式实现相当复杂，本程序通过对问题进行简化有效地实现了\emph{\textbf{基于语义的}}代码相似性判断。

\subsection{两点假设}
为了简化问题，首先进行以下两个假设

\begin{itemize}
	\item 程序的抽象语法树(AST, Abstract Syntax Tree)和实际执行逻辑高度相关
	\item 抽象语法树中的语句节点和表达式节点是所有节点中和实际执行逻辑最相关的两类节点
\end{itemize}

根据这两点假设，通过从程序编译时的抽象语法树的语句(Statement)节点和表达式(Expression)节点序列中寻找最长上升子序列可以有效地衡量程序的逻辑相似性。定义逻辑相似度$s$，其中$AST_j$是程序代码$j$的抽象语法树的语句节点和表达式节点序列

$$
s=\frac{|LCS_{AST_1,AST_2}|}{\mathop{\max}\{|AST_1|,|AST_2|\}}
$$

\clearpage
\appendix

\section{附录： 代码实现}
\label{sec:app_impl}


\begin{lstlisting}[language=Python,numbers=left,style=PythonStyle,caption=矩阵连乘,label={code:implmat}]
p = list([5, 10, 3, 12, 5, 50, 6])
N = 6

m = list([list([0x7fffffff for i in range(0, N + 1)]) for j in range(0, N + 1)])
s = list([list([0x7fffffff for i in range(0, N + 1)]) for j in range(0, N + 1)])


def pretty_print(i, j):
    if i == j:
        print('A{}'.format(i), end=' ')
    else:
        print("(", end=' ')
        pretty_print(i, s[i][j])
        pretty_print(s[i][j] + 1, j)
        print(")", end=' ')


for i in range(1, N + 1):
    m[i][i] = 0

for i in range(N, 0, -1):
    for j in range(i, N + 1):
        if i == j:
            m[i][j] = 0
        else:
            for k in range(i, j):
                val = m[i][k] + m[k + 1][j] + p[i - 1] * p[j] * p[k]
                if m[i][j] > val:
                    m[i][j] = val
                    s[i][j] = k

for i in range(1, N + 1):
    for j in range(1, N + 1):
        print("inf" if m[i][j] == 0x7fffffff else m[i][j], end=' ' if j != N else '\n')

for i in range(1, N + 1):
    for j in range(1, N + 1):
        print("None" if s[i][j] == 0x7fffffff else s[i][j], end=' ' if j != N else '\n')

pretty_print(1, N)

\end{lstlisting}

\end{spacing}
\end{document}